\documentclass[11pt]{article}
\usepackage[margin=1in]{geometry}
\usepackage{mathpazo}
\usepackage{amsmath,amssymb}
\usepackage{microtype}
\usepackage{booktabs}
\usepackage[dvipsnames]{xcolor}
\usepackage[most]{tcolorbox}
\usepackage{enumitem}
\usepackage{hyperref}
\hypersetup{colorlinks=true,linkcolor=MidnightBlue,urlcolor=MidnightBlue}
\urlstyle{tt}
\newcommand{\file}[1]{\nolinkurl{#1}}
\newcommand{\code}[1]{\texttt{#1}}
\newcommand{\mhh}{m_{hh}}
\newcommand{\sbin}{\sigma_\text{bin}}
\newcommand{\stot}{\sigma_\text{tot}}

\definecolor{trainable}{HTML}{1B6CA8}
\definecolor{physics}{HTML}{2E933C}
\definecolor{warn}{HTML}{C0392B}

\newtcolorbox{keyidea}[1][]{enhanced,breakable,colback=physics!6,colframe=physics!70!black,
  fonttitle=\bfseries,title={#1},boxrule=0.6pt,arc=2pt,left=6pt,right=6pt}
\newtcolorbox{aside}[1][]{enhanced,breakable,colback=trainable!5,colframe=trainable!70!black,
  fonttitle=\bfseries,title={#1},boxrule=0.6pt,arc=2pt,left=6pt,right=6pt}
\newtcolorbox{watch}[1][]{enhanced,breakable,colback=warn!5,colframe=warn!70!black,
  fonttitle=\bfseries,title={#1},boxrule=0.6pt,arc=2pt,left=6pt,right=6pt}
\newtcolorbox{assume}[1][]{enhanced,breakable,colback=gray!5,colframe=gray!50!black,
  fonttitle=\bfseries,title={#1},boxrule=0.6pt,arc=2pt,left=6pt,right=6pt}

\setlist{itemsep=2pt,topsep=3pt}

\title{\textbf{Redesigning \texttt{heft\_train\_curriculum.ipynb}:}\\
\textbf{revised plan}\\[4pt]
\large For approval before the notebook is changed\\
\normalsize Supersedes \texttt{docs/curriculum\_plan.pdf} (October 29, 2026)}
\date{October 29, 2026}

\begin{document}
\maketitle

\begin{abstract}
\noindent
This revised plan expands the original three curriculum fixes to include two
additional changes that were decided after the initial plan was written. First,
the custom loss functions (\code{CHI2} and the custom \code{L1Loss}) are
removed in favour of PyTorch's built-in \code{nn.L1Loss}, which has proved
more stable and gives better results. Second, \code{noise\_var} is removed as
a training target: because the data obey Poisson statistics with a known
scale, the variance is analytically determined by the cross section itself and
need not be learned. Together with the three curriculum fixes --- corrected
sort direction, a new \code{DIFFICULTY} switch, and a new \code{PACING} switch
--- the notebook is updated in five targeted places. No cell is modified until
you approve.
\end{abstract}

% ===========================================================================
\section{Changes 1--2: loss function and training target}

\subsection{Revert to \texttt{nn.L1Loss}}

The current notebook offers two loss options, selected by a \code{LOSS}
switch: a custom \code{CHI2} (Pearson $\chi^2$) loss and a custom
\code{L1Loss} with a noise floor. Testing showed that results with the custom
\code{L1Loss} are worse than with the earlier, simpler approach. The fix is to
remove both custom loss classes and replace them with PyTorch's built-in
\code{nn.L1Loss}, which computes the mean absolute error (MAE) between
the model's prediction $\hat{f}_k$ and the true bin cross section $\sigma_k$:
\[
  \mathcal{L} = \frac{1}{K} \sum_{k=1}^{K} \bigl|\hat{f}_k - \sigma_k\bigr|.
\]
This treats every bin equally, with no weighting by cross section or noise
estimate. It is numerically stable and requires no auxiliary inputs.

The \code{LOSS} switch is retained in the switches cell for documentary
clarity, but it now simply instantiates \code{nn.L1Loss()} with no branching.
The custom loss class definitions are removed entirely. A comment in the
switches cell explains why a custom loss was tried and why the built-in is
preferred.

\subsection{Remove \code{noise\_var} from the training target}

The target array currently has two columns,
\[
  Y = \begin{bmatrix} \sigma_k & \hat{v}_k \end{bmatrix},
\]
where $\hat{v}_k$ is a noise-variance estimate included for use in the CHI2
loss denominator. With that loss gone, the second column serves no purpose in
training. The new target array is simply
\[
  Y = \begin{bmatrix} \sigma_k \end{bmatrix}.
\]

\begin{aside}[Why \texttt{noise\_var} is not needed as a target]
The data are scaled Poisson counts. Each spectrum $j$ consists of $N_\text{ev}
= 40\,000$ events with equal weight $w_j$; bin $k$ receives $n_k$ events, so
\[
  \sigma_k = w_j\,n_k, \qquad
  \operatorname{Var}(\sigma_k) = w_j^2\,n_k = w_j\,\sigma_k.
\]
The variance is therefore determined exactly by the cross section itself, up to
the per-spectrum weight $w_j$. There is no free parameter to learn: any time
a variance estimate is needed (for diagnostics, for a future weighted loss, or
for uncertainty propagation), it can be computed analytically as $w_j\,
\hat{\sigma}_k$ from the model's own prediction $\hat{\sigma}_k$.
\end{aside}

The \code{noise\_var} column is \emph{kept in the dataframe}; removing it from
training does not require removing it from the data file. This preserves it
for diagnostic use.

\paragraph{Cells that change.} Two cells are affected by these two changes.

\begin{enumerate}
  \item \textit{Switches cell:} \code{LOSS} is set to \code{'L1'} and the
    custom loss class definitions are removed. A one-line comment explains what
    was tried and why the built-in is used instead.
  \item \textit{Y-array construction cell:} the line
    \begin{quote}\small\ttfamily
      Y = df[[NN.TARGET, 'noise\_var']].to\_numpy(dtype=np.float32)
    \end{quote}
    becomes
    \begin{quote}\small\ttfamily
      Y = df[[NN.TARGET]].to\_numpy(dtype=np.float32)
    \end{quote}
    and the training loop is updated to pass \code{y.squeeze()} (or
    \code{y[:,\,0]}) to the loss function rather than unpacking two columns.
\end{enumerate}

% ===========================================================================
\section{Changes 3--5: the curriculum}

These three changes are unchanged from the original plan. They are reproduced
here for completeness.

\subsection*{Background: easy means large cross section}

Bengio et al.\ (2009) showed that training on examples ordered from easy to
hard often converges faster and to a better minimum. For \textsc{HEFTNet}, a
training row is one bin $(x_k,\,\theta,\,\sigma_k)$. The SNR of that bin is
\[
  \mathrm{SNR}_k = \sqrt{\frac{\sigma_k}{w_j}},
\]
so a large $\sigma_k$ means a reliable gradient and a small $\sigma_k$ means
a noisy one. \textbf{Easy rows are those with the largest cross section, not
the smallest.}

\subsection{Fix the sort direction (Change 3)}

Cell~19 currently sorts in ascending order:
\begin{quote}\small\ttfamily
order = np.argsort(key, kind='stable')
\end{quote}
This puts the hardest (smallest) examples first. The fix is:
\begin{quote}\small\ttfamily
order = np.argsort(score, kind='stable')[::-1]
\end{quote}
The variable is renamed \code{score} for clarity, and a comment explains why
descending order is correct.

\subsection{Replace \texttt{CURRICULUM\_BY} with \texttt{DIFFICULTY} (Change 4)}

The existing switch conflates the difficulty criterion (what makes an example
easy) with the level at which it is applied (per-bin or per-spectrum). These
are split into a single, clearly named switch:

\begin{center}
\small
\begin{tabular}{lll}
\toprule
\code{DIFFICULTY} value & sort key & granularity \\
\midrule
\code{'bin\_sigma'} & $\sigma_k$ of that bin & per bin \\
\code{'spectrum\_sigma'} & $\sigma_\text{tot} = \sum_k \sigma_k$ of the parent
  spectrum & per spectrum \\
\code{'spectrum\_snr'} & $\sigma_\text{tot} / \sqrt{w_j}$ & per spectrum \\
\bottomrule
\end{tabular}
\end{center}

The \code{'spectrum\_snr'} option requires the per-spectrum event weight $w_j$.
This is not in the current training file; the option is included in the switch
for future use but raises \code{NotImplementedError} if selected, with a
comment explaining what to add.

\subsection{Add \texttt{PACING} switch and log ramp (Change 5)}

A new \code{PACING} switch controls the shape of the fraction ramp:

\begin{itemize}
  \item \code{'linear'}: $f(i) = \min\bigl(1,\; f_0 + (1-f_0)\,i/i_\text{ramp}\bigr)$.
    This is the current behaviour.
  \item \code{'log'}: $f(i) = \min\bigl(1,\; f_0^{\,1 - i/i_\text{ramp}}\bigr)$.
    At $i=0$ this gives $f_0$; at $i = i_\text{ramp}$ it gives 1. The growth
    is slow at first and accelerates, spending more iterations on the easiest
    examples before the ramp opens up.
\end{itemize}

The \code{CurriculumLoader.fraction} method is updated with an \code{if/else}
block for the two cases, and the constructor takes \code{pacing} as an
argument. The docstring explains both shapes and why the sort order matters.

The \code{DIRNAME} string is updated from \code{\_curriculum\_bin} /
\code{\_curriculum\_spectrum} to \code{\_difficulty\_<value>\_pacing\_<value>}
so that run folders are self-describing.

% ===========================================================================
\section{Summary of all five changes}

\begin{center}
\small
\begin{tabular}{cllll}
\toprule
\# & cell & what changes & replaces & status \\
\midrule
1 & Switches & remove custom loss classes; use \code{nn.L1Loss} &
  custom \code{CHI2} / \code{L1Loss} & new \\
2 & Y-array & drop \code{'noise\_var'} column from \code{Y} &
  two-column target & new \\
3 & Sorting & add \code{[::-1]} for descending order &
  ascending sort & from v1 \\
4 & Switches & \code{DIFFICULTY} replaces \code{CURRICULUM\_BY} &
  \code{CURRICULUM\_BY} & from v1 \\
5 & \code{CurriculumLoader} & add \code{PACING} and log ramp &
  linear only & from v1 \\
\bottomrule
\end{tabular}
\end{center}

% ===========================================================================
\section{Experiments to run after the changes}

Four short runs (3$\times$20k iterations, seeded, \code{nn.L1Loss}):

\begin{center}
\small
\begin{tabular}{llll}
\toprule
run & curriculum & \code{DIFFICULTY} & \code{PACING} \\
\midrule
baseline & no & --- & --- \\
A & yes & \code{'bin\_sigma'} & \code{'linear'} \\
B & yes & \code{'bin\_sigma'} & \code{'log'} \\
C & yes & \code{'spectrum\_sigma'} & \code{'log'} \\
\bottomrule
\end{tabular}
\end{center}

Three seeds per run give enough spread to assess significance. The baseline
run also uses \code{nn.L1Loss}, so the comparison is clean. Once a best
curriculum setting is identified, a full 200k-iteration run follows.

\begin{watch}[Previous curriculum results are no longer comparable]
The curriculum runs recorded in the report used the wrong sort direction
\emph{and} a custom loss function. They are not a valid baseline for the new
experiments. All comparisons should be made against the freshly seeded
\code{nn.L1Loss} baseline above.
\end{watch}

% ===========================================================================
\section{What counts as success}

A curriculum is worth keeping if, compared to the plain baseline with the same
seeds and same number of iterations, it achieves a statistically meaningful
improvement on the 1\,000 holdout spectra: a lower median $\chi^2/N$
(evaluated \emph{offline}, not used for training) by more than the
run-to-run spread, and/or a lower fraction of spectra with $\chi^2/N > 2$.
If neither criterion is met, the plain \code{nn.L1Loss} training is the
default.

% ===========================================================================
\section{Changes not made by this plan}

\begin{itemize}
  \item \code{DIFFICULTY = 'spectrum\_snr'} is stubbed out (\code{NotImplementedError});
    it needs the per-spectrum event weight $w_j$ added to the training file.
  \item Seeding (\code{torch.manual\_seed} and \code{numpy.random.seed}) is a
    separate one-cell change, to be applied to both \file{heft\_train.ipynb}
    and \file{heft\_train\_curriculum.ipynb} together.
  \item Nothing changes in \file{heft\_test.ipynb}, \file{heft\_checks.ipynb},
    or \file{heft\_export.ipynb}.
  \item The \code{noise\_var} column is kept in the dataframe and data files;
    only its inclusion in the training target array is removed.
\end{itemize}

% ===========================================================================
\begin{keyidea}[Decisions for you before I proceed]
\begin{enumerate}
  \item \textbf{Loss function.} Confirmed: use \code{nn.L1Loss} and remove
    the custom loss classes. Any objection?
  \item \textbf{Noise-var target.} Confirmed: remove \code{'noise\_var'} from
    \code{Y}. Keep the column in the dataframe. Any objection?
  \item \textbf{Sort direction.} Easy $=$ large cross section $=$ descending
    sort. Any objection?
  \item \textbf{Difficulty options.} \code{'bin\_sigma'},
    \code{'spectrum\_sigma'}, \code{'spectrum\_snr'} (stubbed). Any additions
    or changes?
  \item \textbf{DIRNAME naming.}
    \file{runs/L1Loss\_difficulty\_bin\_sigma\_pacing\_log}. Acceptable, or
    should we abbreviate?
  \item \textbf{Seeding.} Same time as these changes, or separately
    afterwards?
\end{enumerate}
\end{keyidea}

\end{document}
